% Part: incompleteness
% Chapter: arithmetization-syntax
% Section: coding-symbols

\documentclass[../../../include/open-logic-section]{subfiles}

\begin{document}

\olfileid{inc}{art}{cod}
\olsection{Pangodean Simbol}

Basa dhasar~$\Lang L$ logika tataran kapisan
nggunakake simbol-simbol
\[
\lfalse \quad \lnot \quad \lor \quad \land \quad \lif \quad \lforall
\quad \lexists \quad \eq \quad ( \quad ) \quad ,
\]
bebarengan karo himpunan !!{enumerable}
variabel lan !!{constant}s, sarta himpunan
!!{enumerable} !!{function}s lan
!!{predicate}s kanthi aritas sembarang.
Kita bisa menehi \emph{kode} marang saben
simbol iki saengga saben simbol nduwé siji
wilangan khas minangka kodene, lan ora ana
rong simbol béda kang diwenehi wilangan
kang padha. Iki mungkin amarga himpunan
kabeh simbol iku !!{enumerable}, mula ana
!!a{bijection} antarane himpunan iku lan
himpunan wilangan natural. Nanging kita uga
kepengin bisa nemokake maneh simbol kasebut
(bebarengan karo sawatara katrangan, umpamane
aritas !!a{function}) saka kodene kanthi
cara komputabel. Mesthi ana akeh cara kanggo
nindakake iku. Ing ngisor iki salah siji cara
kang nggunakake fungsi rekursif primitif.
(Elinga yen $\tuple{n_0, \dots, n_k}$ iku
wilangan kang ngodeake urutan wilangan
$n_0$, \dots, $n_k$.)

\begin{defn}
Yen $s$ iku simbol saka~$\Lang L$,
tetepna \emph{kode simbol}~$\scode s$
kaya mangkene:
\begin{enumerate}
\item Yen $s$ kalebu simbol logis,
  $\scode s$ diwenehi dening tabel iki:
\[
\begin{array}{cccccccccc}
  \lfalse & \lnot & \lor & \land & \lif & \lforall \\
\tuple{0, 0} & \tuple{0, 1} & \tuple{0, 2} & \tuple{0, 3} &
\tuple{0, 4} & \tuple{0, 5} \\
\lexists & \eq & ( & ) & ,\\
\tuple{0, 6} & \tuple{0, 7} &
\tuple{0, 8} & \tuple{0, 9} & \tuple{0, 10}
\end{array}
\]
\item Yen $s$ iku variabel kapangka-$i$
  $\Obj v_i$, mula $\scode s = \tuple{1, i}$.
\item Yen $s$ iku !!{constant} kapangka-$i$
  $\Obj c_i$, mula $\scode s = \tuple{2, i}$.
\item Yen $s$ iku !!{function} kapangka-$i$
  kanthi aritas~$n$ $\Obj f_i^n$, mula
  $\scode s = \tuple{3, n, i}$.
\item Yen $s$ iku !!{predicate} kapangka-$i$
  kanthi aritas~$n$ $\Obj P_i^n$, mula
  $\scode s = \tuple{4, n, i}$.
\end{enumerate}
\end{defn}

\begin{prop}
Relasi-relasi ing ngisor iki rekursif primitif:
\begin{enumerate}
\item $\fn{Fn}(x, n)$ yen lan mung yen $x$
  iku kode $\Obj f^n_i$ kanggo sawijining~$i$,
  yaiku $x$ iku kode !!{function} aritas~$n$.
\item $\fn{Pred}(x, n)$ yen lan mung yen $x$
  iku kode $\Obj P^n_i$ kanggo sawijining~$i$
  utawa $x$ iku kode $\eq$ lan $n = 2$,
  yaiku $x$ iku kode !!{predicate} aritas~$n$.
\end{enumerate}
\end{prop}

\begin{defn}
Yen $s_0, \dots, s_{n-1}$ iku urutan simbol,
\emph{wilangan G\"odel} saka urutan iku yaiku
$\tuple{\scode{s_0}, \dots, \scode{s_{n-1}}}$.
\end{defn}

\begin{explain}
Gatekna yen \emph{kode} lan \emph{wilangan G\"odel}
iku rong bab kang béda. Tuladhane,
variabel~$\Obj v_5$ nduwé kode~$\scode{\Obj
  v_5} = \tuple{1, 5} = 2^2\cdot 3^6$.
Nanging variabel~$\Obj v_5$ yen dianggep
minangka suku uga dadi urutan simbol
(dawane~$1$). \emph{Wilangan G\"odel}~$\Gn{\Obj v_5}$
saka \emph{suku}~$\Obj v_5$ yaiku
$\tuple{\scode{\Obj v_5}} = 2^{\scode{\Obj
    v_5} + 1} = 2^{2^2\cdot 3^6 + 1}$.
\end{explain}

\begin{ex}
Elinga yen $k_0$, \dots, $k_{n-1}$ iku
urutan wilangan. Kode urutan
$\tuple{k_0, \dots, k_{n-1}}$ ing
pangodean pangkat wilangan prima yaiku
\[
2^{k_0+1}\cdot3^{k_1+1}\cdot \dots \cdot p_{n-1}^{k_{n-1}+1},
\]
ing ngendi $p_i$ iku wilangan prima kapangka-$i$
(diwiwiti saka $p_0 = 2$). Mula, umpamane,
formula $\eq[\Obj v_0][\Obj 0]$, utawa
kanthi luwih cetha, ${\eq}(\Obj v_0,\Obj c_0)$,
nduwé wilangan G\"odel
\[
\tuple{\scode{\eq},\scode{(},\scode{\Obj v_0},\scode{,}, \scode{\Obj
    c_0},\scode{)}}.
\]
Ing kene, $\scode{\eq}$ yaiku
$\tuple{0,7} = 2^{0+1}\cdot 3^{7+1}$,
$\scode{\Obj v_0}$ yaiku
$\tuple{1,0} = 2^{1+1}\cdot3^{0+1}$,
lan sateruse. Mula $\Gn{=(\Obj v_0,\Obj c_0)}$ yaiku
\begin{multline*}
2^{\scode{=} + 1}\cdot 3^{\scode{(}+1}\cdot 5^{\scode{\Obj v_0}+1}
\cdot 7^{\scode{,} + 1} \cdot 11^{\scode{\Obj c_0}+1} \cdot
13^{\scode{)}+1} = \\
2^{2^1\cdot 3^8 + 1}\cdot 3^{2^1\cdot 3^9+1}\cdot 5^{2^2\cdot 3^1+1}
\cdot 7^{2^1\cdot 3^{11} + 1} \cdot 11^{2^3\cdot3^1+1} \cdot
13^{2^1\cdot3^{10}+1} = \\
2^{13\,123}\cdot 3^{39\,367}\cdot 5^{13}\cdot 7^{354\,295}\cdot11^{25}\cdot13^{118\,099}.
\end{multline*}
\end{ex}

\end{document}
